
\documentclass[a4paper,12pt]{article}

\input{glyphtounicode}
\pdfgentounicode=1

\usepackage[T1]{fontenc}
\usepackage[utf8]{inputenc}
\usepackage[english]{babel}
\usepackage[version=4]{mhchem}
\usepackage[dvipsnames]{xcolor}
\usepackage[backend=biber,style=ieee]{biblatex}
\usepackage{newtxtext,newtxmath}
\usepackage{microtype}
\usepackage{setspace}
\usepackage{titlesec}
\usepackage{tocloft}
\usepackage{indentfirst}
\usepackage{enumitem}
\usepackage{fancyhdr}
\usepackage{hyperref}
\usepackage{tabularx}
\usepackage{array}
\usepackage{makecell}
\usepackage{booktabs}
\usepackage{graphicx}
\usepackage{amsmath}
\usepackage{footmisc}
\usepackage{csquotes}
\usepackage{fontawesome5}
\usepackage[
  a4paper,
  top        = 2.00cm,
  bottom     = 2.00cm,
  left       = 2.00cm,
  right      = 2.00cm,
  headheight = 15pt,
  headsep    = 5pt,
  footskip   = 15pt
]{geometry}

\newcommand{\degreeawardvalue}{}
\newcommand{\universityvalue}{}
\newcommand{\addressvalue}{}
\newcommand{\unilogovalue}{}
\newcommand{\copyyearvalue}{}
\newcommand{\defenddatevalue}{}
\newcommand{\orcidvalue}{}
\newcommand{\rightsstatementvalue}{}
\newcommand{\degreeaward}[1]{\def\degreeawardvalue{#1}}
\newcommand{\university}[1]{\def\universityvalue{#1}}
\newcommand{\address}[1]{\def\addressvalue{#1}}
\newcommand{\unilogo}[1]{\def\unilogovalue{#1}}
\newcommand{\copyyear}[1]{\def\copyyearvalue{#1}}
\newcommand{\defenddate}[1]{\def\defenddatevalue{#1}}
\newcommand{\orcid}[1]{\def\orcidvalue{#1}}
\newcommand{\rightsstatement}[1]{\def\rightsstatementvalue{#1}}
\newcommand{\titleskip}{\vskip 4\bigskipamount}
\newcommand{\authorskip}{\vskip 2\bigskipamount}
\newcommand{\resumesection}[2]{\section{\textcolor{red}{#1}#2}}
\newcommand{\resumeSubHeadingListStart}{\begin{itemize}[leftmargin=0.0in, label={}, itemsep=2pt]}
\newcommand{\resumeSubHeadingListEnd}{\end{itemize}}
\newcommand{\resumeItemListStart}{\begin{itemize}}
\newcommand{\resumeItemListEnd}{\end{itemize}}
\newcommand{\resumeItem}[1]{
    \item\small{#1}
}
\newcommand{\resumeSubheading}[4]{
    \item
    \begin{tabular*}{1.0\textwidth}[t]{l@{\extracolsep{\fill}}r}
        \textbf{#1} & \textbf{\small #2} \\
        \small#3 & \small #4
    \end{tabular*}
}

\renewcommand{\cftsecleader}{\cftdotfill{\cftdotsep}}
\renewcommand{\cftsecpresnum}{Chapter~}
\renewcommand{\cftsecaftersnum}{:}
\renewcommand{\headrulewidth}{0pt}
\renewcommand{\footrulewidth}{0pt}
\renewcommand\labelitemi{$\vcenter{\hbox{\tiny$\bullet$}}$}
\renewcommand\labelitemii{$\vcenter{\hbox{\tiny$\bullet$}}$}

\pagestyle{fancy}
\fancyhf{}
\fancyfoot[R]{\thepage}
\urlstyle{same}
\addbibresource{references.bib}
\hypersetup{
  colorlinks=true,
  linkcolor=RoyalBlue,
  urlcolor=RoyalBlue,
  citecolor=RoyalBlue,
  anchorcolor=RoyalBlue
}
\settowidth{\cftsecnumwidth}{Chapter 9: }
\titleformat{\section}
  {\normalfont\Large\bfseries}
  {Chapter \thesection:}
  {0.5em}
  {}
% \titleformat{\section}{
%     \scshape\raggedright\large\bfseries
%     }{}{0em}{}[\color{black}\titlerule]
%     \titlespacing{\section}{0pt}{5pt}{5pt}

\setstretch{1.5}
\setlist[itemize]{topsep=4pt, itemsep=3pt, parsep=0pt, partopsep=0pt, leftmargin=*}
\setlist[enumerate]{topsep=4pt, itemsep=3pt, parsep=0pt, partopsep=0pt, leftmargin=*}
\setlength{\footnotesep}{15pt} 
\setlength{\skip\footins}{15pt} 
\setlength{\parskip}{5pt}

\begin{document}

{\centering
  {\scshape \textls[50]{\fontsize{25}{30}\selectfont Reza {\bfseries Aliasgari Renani}}} \\
  {\fontsize{10}{10}\selectfont 11/1/2026} $|$
  {\fontsize{10}{10}\selectfont Motivation Letter} $|$
  {\fontsize{10}{10}\selectfont Ventspils University of Applied Sciences} $|$
  {\fontsize{10}{10}\selectfont PhD E-Sailors 9} $|$
  {\fontsize{10}{10}\selectfont Lunar Orbit Communications}
\par}
\vspace{-15pt}
\noindent\rule{\textwidth}{0.5pt}

\vspace{-5pt}

My research objective is to contribute to lunar-orbit communications and RF ranging for small-satellite missions by developing robust ground-segment experiments and data-processing workflows. The E-Sailor~9 doctoral position at Ventspils University of Applied Sciences focuses on WP4, ``Lunar Orbit Communications and Ranging: RT16 and RT32 Ground Segment Design,'' with access to the Irbene radio astronomy and space communication complex, including RT-32, RT-16, and the LOFAR station. This scope matches my background in experimental systems, electronics, and data processing. I am completing the final year of my Master's degree in Plasma Physics at Moscow Institute of Physics and Technology, and my experience with mission-relevant hardware and measurement systems has led me to pursue doctoral research in space communications for ESTCube-LuNa.

I joined Dr.\ Koveshnikov's laboratory in the Russian Academy of Sciences over two years ago to study irradiated semiconductor devices. We investigated the effects of electron and gamma irradiation and hydrogen plasma treatments on \ce{SiO2}-based MOS devices, distinguishing trap contributions by location (oxide, semiconductor, interface) and carrier type (electrons, holes) using electrical characterization techniques. This work resulted in a first-author publication\footnote{R.~Aliasgari~Renani, O.A.~Soltanovich, M.A.~Knyazev, S.V.~Koveshnikov, ``Investigation of low energy electron irradiated \ce{SiO2} based MOS devices by C--V and TSC techniques,'' \textit{Russian Microelectronics}, 2023, DOI: \href{https://doi.org/10.1134/S1063739723600516}{10.1134/S1063739723600516}} and strengthened my understanding of device stability under bias stress and radiation. I also developed optimized MATLAB automation programs for experimental techniques such as TSC, C--V, I--V, I--T, and DLTS measurements, enabling systematic data acquisition and reduced experimental run-times.

After my undergraduate degree, I joined the System-on-Chip Development Laboratory under the supervision of Dr.\ Vasiliev to develop FPGA-based systems. I implemented video-processing algorithms in Verilog using both manual coding and HDL generated from Simulink models, built verification testbenches, and used Python for quality control. I developed a fixed-point arithmetic library and LUT-based function approximations using Horner's method. I evaluated resource utilization (LUTs and BRAM) and synthesized the designs on a Xilinx Zybo Z7 board with live camera input and HDMI output. The codebase was maintained under Git-based source control. Concurrently, at the Laboratory of Plasma Systems, I conducted research on radiation and plasma effects on FPGA devices. We exposed an FPGA board to 25--60~keV electron beams in low-pressure oxygen, generating plasma and X-rays from a tungsten target while applying thermal cycling. The FPGA output signal was continuously monitored, and preliminary functionality tests were performed during and after exposure.

Within the E-Sailor~9 PhD project, I aim to contribute to requirements and specification work for ESTCube-LuNa communications via the RT-16 and RT-32 ground stations, design and validate the ground-segment ranging experiment, and develop signal-processing and simulation workflows. This includes building and documenting experimental setups, implementing data-processing pipelines, and validating RF ranging and communication chains through reproducible analyses and test campaigns, including preliminary experiments using other spacecraft as signal sources. I intend to document results in technical reports and peer-reviewed publications and to contribute to the broader E-Sailors collaboration through workshops and conferences.

This doctoral position represents a continuation of my experience working with experimental hardware, measurement systems, and data analysis in space-relevant contexts. The combination of doctoral training at Ventspils University of Applied Sciences, access to the Irbene radio astronomy and space communication complex, and collaboration with KTH provides an environment well suited to developing expertise in lunar-orbit communications and ground-segment experimentation. I am particularly interested in working under the supervision of Prof.\ Andris Vaivads and collaborating with Prof.\ Mykola Ivchenko within the E-Sailors framework.

\end{document}
